\documentclass[10pt,a4paper]{amsart}
\usepackage[margin=19mm]{geometry}
\usepackage{amsmath,amssymb,mathtools}
\usepackage[scr=rsfs,cal=euler]{mathalfa}
\usepackage{xfrac}
\usepackage[unicode,hidelinks,bookmarks=false]{hyperref}
\newcommand{\ie}{i.e.\ }
\newcommand{\cf}{cf.\ }
\newcommand{\resp}{respectively\ }
\newcommand{\Subsection}{\subsection}
\providecommand{\nobreakdash}{\nobreak\hbox{-}}
\setlength{\emergencystretch}{2em}
\multlinegap0pt
\usepackage{bm}
\usepackage{bbm}
\usepackage{dsfont}
\usepackage{xspace}
\usepackage{cancel}
\usepackage{mathtools}
\usepackage{amsthm}
\makeatletter
\@ifundefined{lemma}{\newtheorem{lemma}{Lemma}}{}
\makeatother
\title[A non-sequential arithmetical theory with pairing]{A non-sequential arithmetical theory with pairing}
\author{Juvenal Murwanashyaka}
\date{Source: arXiv:2509.15191v1 (CC BY 4.0)}
\begin{document}
\maketitle

\noindent\emph{Source and licence.} A non-sequential arithmetical theory with pairing. Version arXiv:2509.15191v1 (CC BY 4.0), distributed under the Creative Commons Attribution 4.0 International licence, as recorded in the arXiv licence metadata for this exact version and shown on the abstract page https://arxiv.org/abs/2509.15191v1. Full rights notice: licenses/arxiv-2509.15191-CC-BY-4.0.txt.\par\smallskip

\noindent\textit{Editorial scope.} Counted unit: the Lemma whose source label is \texttt{SectionQNthClureLemma} in the article (source0.tex lines 2122--2129, source label \texttt{SectionQNthClureLemma}), delivered together with the article's own complete original proof of it (source0.tex lines 2130--2521, one \texttt{proof} environment of 392 lines). No mathematics is written, repaired or re-derived: statement and proof are the article's own text line for line. Required context retained uncounted: none. Every definition, display and label of those lines is retained unchanged. Omitted from this package and not redistributed: 1--2121, 2522--4154. Nothing inside a retained excerpt is omitted or altered: every retained body is delivered exactly as the article's lines are written, so no omission receipt of any kind accompanies this package. Citations: the retained excerpts carry no citation command at all, so no bibliography entry of the article is transcribed and no reference is redistributed. Reference adaptation: none was needed and none was made; every reference inside the delivered unit resolves inside it, so no unresolved reference or citation is produced. Printed slips kept literal: no printed word, symbol, hypothesis or number of the counted statement or of its proof was altered. Layout and class: the article's own class is \texttt{memoir} with packages including inputenc, babel, amsmath, amsthm, thmtools, mathtools, amsfonts, amssymb, makeidx, graphicx, tikz-cd, float, authblk, forest; the wrapper is the lane's pinned \texttt{amsart} editorial wrapper at 10pt on A4, which restates the theorem-like environments the retained text needs and adds the article's own macro definitions that the retained text needs (none), with extra wrapper packages (bm, bbm, dsfont, xspace, cancel, mathtools). The delivered numbering is the unit's own. Assets: no retained span contains a figure environment, a graphics-inclusion command, a PDF-page inclusion, a table or a TikZ picture; no image file is copied into this package, the package declares no asset dependency and no image file is redistributed with it. Licence: version arXiv:2509.15191v1 of this article is distributed under the Creative Commons Attribution 4.0 International licence, as recorded in the arXiv licence metadata for this exact version and shown on the abstract page https://arxiv.org/abs/2509.15191v1. A full rights notice is delivered at licenses/arxiv-2509.15191-CC-BY-4.0.txt. Characters: every retained excerpt is pure ASCII, so the delivered engine, XeLaTeX with its default Latin Modern text font, reports no missing character.
\begin{lemma}   \label{SectionQNthClureLemma}
Let  $X \subseteq   \mathbb{T} $  and   $  k, N  \in  \omega $. 
Assume    $  F :  \mathcal{W}_{ k + N  }  \left(  X  \right)   \to   \mathbb{T}    $  is  a  \textsf{Good}  $  L_{  \mathbb{T} }^- $-embedding.
Then,  there  exists  a   one-to-one  $  L_{  \mathbb{T} }^-$-homomorphism   
 $  L:   \mathsf{Cl}_{N}  \left(  \mathcal{W}_{k}  \left( X  ,  \mathsf{d}^0_0 \right)   \right)   \to   \mathbb{T}   $
that agrees with  $  F  $  on     $  \mathcal{W}_{ k   }  \left(  X  \right)    $
and  fixes    $  \big\{   \mathsf{d}^m_0  :   \    m  \in  \mathbb{Z}   \,   \big\}     \,    $.
\end{lemma}

\begin{proof}


Observe that 
\[
\mathcal{W}_{k}  \left( X  ,  \mathsf{d}^0_0 \right)    
= 
  \mathcal{W}_{k}  \left( X   \right)    \cup     \mathcal{W}_{k}  \left(  \mathsf{d}^0_0   \right)  
= 
  \mathcal{W}_{k}  \left( X   \right)  
\cup   \big\{   \mathsf{d}^m_0  :   \    m  \in  \mathbb{Z}   \,   \big\}
\        .
\]
First,  we  introduce a rank function on   $ \mathsf{Cl}_{N}  \left(  \mathcal{W}_{k}  \left( X ,  \mathsf{d}^0_0   \right)   \right)    $
\[
\mathsf{rk}  \left(  w  \right)   =  \min  
\big\{  m  \in  \lbrace 0, 1,  \ldots ,  N  \rbrace  :   \  
 w  \in  \mathsf{Cl}_{m}  \left(  \mathcal{W}_{k}  \left( X ,   \mathsf{d}^0_0    \right)   \right)     \big\}  
\     .
\]
We define  $L$ by recursion:   for  $  w  \in      \mathsf{Cl}_{N}  \left(  \mathcal{W}_{k}  \left( X  ,   \mathsf{d}^0_0    \right)   \right)     $
\[
L  \left(  w  \right)   =   \begin{cases}
w            &   \mbox{  if   }     w  \in    \big\{   \mathsf{d}^m_0  :   \    m  \in  \mathbb{Z}   \,   \big\} 
\\
F ( w)      &     \mbox{  if   }   w  \in    \mathcal{V}_{ k  }  \left(  X  \right) 
\\
\langle  L(w_0),   L(w_0 )   \rangle^{   \mathcal{M}  }      &    \mbox{  if  }  w  \not\in   \mathcal{V}_{ k  }  \left(  X  \right)   
                     \mbox{  and   }    w  =  \langle  w_0,  w_1    \rangle^{   \mathcal{M}  }  
\\
\mathcal{R}_m  \left(   L(v)  \right)      &    \mbox{  if  }  w  \not\in   \mathcal{V}_{ k  }  \left(  X  \right)   
                     \mbox{  and   }      w  =       \mathcal{R}_m  \left(   v \right)    
\\
\mathcal{A}_n  \left(   L(v)  \right)      &    \mbox{  if  }  w  \not\in   \mathcal{V}_{ k  }  \left(  X  \right)   
                     \mbox{  and   }      w  =       \mathcal{A}_n  \left(   v \right)                                   
\end{cases}
\]
where   $   w_0,  w_1,  v   \in   \mathsf{Cl}_{N}  \left(  \mathcal{W}_{k}  \left( X   ,  \mathsf{d}^0_0   \right)   \right)     $, 
$  m  \in  \mathbb{Z}  $  and  $  n \in  \omega  \setminus  \lbrace 0, 1  \rbrace  $. 
This map is well-defined since: 
(i)   $   \mathcal{R}_m  \left(   v \right)    =  \langle    \mathcal{R}_{m-1}  \left(   v \right)     \,  ,   \,  v  \rangle^{  \mathcal{B} }  $  
and  
 $   \mathcal{A}_v  \left(   v \right)    =  \langle    \mathcal{A}_{n-1}  \left(   v \right)     \,  ,   \,  v  \rangle^{  \mathcal{B} }  $; 
 (ii) distinct maps in    $   \lbrace  \mathcal{R}_m   :  \   m  \in  \mathbb{Z} \,   \rbrace   \cup 
 \lbrace  \mathcal{A}_n :  \   n  \in    \omega  \setminus \lbrace 0, 1   \rbrace  \,   \rbrace       $
 have disjoint images; 
 (iii)  $F$ is  \textsf{Good}.
Recall that    $  \mathcal{A}_1  $  is  the identity map. 










To show that $L$ is one-to-one,  we need the following claim. 

   
   \begin{quote} {\bf (Claim)} \;\;\;\;\;\;    
   Let    $   v  \in    \mathsf{Cl}_{N}  \left(  \mathcal{W}_{k}  \left( X  \right)     ,  \mathsf{d}^0_0     \right)   $     and 
   $  w  \in  \mathcal{W}_{k + N- \mathsf{rk}  \left( v  \right)  }  \left( X  \right)    $
   be distinct. 
Then,    $ L(v)   \neq   F(w) $.
 \end{quote}  
 
 Pick   $   v  \in    \mathsf{Cl}_{N}  \left(  \mathcal{W}_{k}  \left( X    ,  \mathsf{d}^0_0    \right)   \right)   $. 
 We prove  by induction on  $  \mathsf{rk}   \left(  v  \right)  $  that  
 $  L(v)  \neq  F(w)  $  for all   $  w    \in \mathcal{W}_{k + N- \mathsf{rk}  \left( v  \right)  }  \left( X  \right)   \setminus  \lbrace  v  \rbrace  $. 
 We consider the base case   $   \mathsf{rk}_N  \left(  v  \right)  = 0  $. 
 We have two cases:   
 (1)  $  v    \in  \mathcal{W}_{k}  \left( X  \right)     $; 
 (2)  $  v  =  \mathsf{d}^m_0  $  where  $  m  \in  \mathbb{Z}  $. 
 In case of  (1),  since   $L$ and $F$ agree on   $   \mathcal{W}_{k}  \left( X  \right)     $ and  $F$  is one-to-one, 
  $  L(v)   =  F(v)   \neq    F(w)  $    for all   $  w    \in \mathcal{W}_{k + N }  \left( X  \right)   \setminus  \lbrace  v  \rbrace  $. 
   We consider (2):   $  v  =  \mathsf{d}^m_0     \,  $.
   We have  $  L(\mathsf{d}^m_0 )  =  \mathsf{d}^m_0   $. 
   Pick    $  w    \in \mathcal{W}_{k + N }  \left( X  \right)   \setminus  \lbrace  v  \rbrace  $. 
   Assume for the sake of a contradiction  $ \mathsf{d}^m_0   =  F(w)  $. 
   Since  $  w  \neq  \mathsf{d}^m_0   =  F(w)  $  and  $  F $  is  \textsf{Good},  
   we must have  $  w,  F(w)  \not\in    \big\{   \mathsf{d}^m_0  :   \    m  \in  \mathbb{Z}   \,   \big\}  $, 
   which contradicts the assumption that    $  \mathsf{d}^m_0   =  F(w)  $. 
   Thus,     $  L(v)   \neq    F(w)  $    for all   $  w    \in \mathcal{W}_{k + N }  \left( X  \right)   \setminus  \lbrace  v  \rbrace  $. 

   
   
   
   
   
We consider the case      $   \mathsf{rk}  \left(  v  \right)   >  0  $. 
We   have the following cases: 
 \begin{enumerate}
\item  $  v  =  \langle  s,  t  \rangle^{  \mathcal{B}  }  $  where  
$  s,  t  \in    \ \mathsf{Cl}_{N}  \left(  \mathcal{W}_{k}  \left( X   ,   \mathsf{d}^0_0    \right)   \right)   $
and  $  \mathsf{rk} \left(  v  \right)  =  1  +  \max\left(   \mathsf{rk} \left(  s \right)  \,  ,  \,    \mathsf{rk} \left(  t  \right)   \right)    \,   $. 




\item  $  v  =   \mathcal{R}_m \left(  s  \right)   $  where  
$  s  \in     \mathsf{Cl}_{N}  \left(  \mathcal{W}_{k}  \left( X    ,   \mathsf{d}^0_0    \right)   \right)   $, 
  $  m  \in  \mathbb{Z}  $
  and  $  \mathsf{rk} \left(  v  \right)  =  1  +  \mathsf{rk} \left(  s \right)      \,   $. 



 
 
 
 

\item  $  v  =   \mathcal{A}_n \left(  s  \right)   $  where  
$  s  \in    \mathsf{Cl}_{N}  \left(  \mathcal{W}_{k}  \left( X     ,   \mathsf{d}^0_0      \right)   \right)    $, 
  $ n  \in  \omega  \setminus  \lbrace 0, 1  \rbrace   $
  and  $  \mathsf{rk} \left(  v  \right)  =  1  +  \mathsf{rk} \left(  s \right)      \,   $. 
 \end{enumerate}
Pick     $  w  \in  \mathcal{W}_{k + N- \mathsf{rk} \left( v  \right)  }  \left( X  \right)    $
such that  $  L(v)  =  F(w)  $.  
We need to show that  $  v  =  w  $. 
We consider  case  (1). 
Since  $L$ is an $L_{  \mathbb{T} }^-  $-homomorphism  and  $F$  is  an   $L_{  \mathbb{T} }^-  $-embedding
\[
\langle  L(s), L(t)   \rangle^{  \mathcal{B}  } 
= L(v)  = F(w)   \in  
 \mathcal{W}_{k + N- \mathsf{rk} \left(   v \right)  }  \left(   F  \left([ X  \right]   \right)   
 \     . 
\]
Since  $   k + N- \mathsf{rk}  \left(   v \right)  +1  \leq  k+N $  as  $   \mathsf{rk}  \left(   v \right)   >  0  $
\[
  L(s), L(t)  \in       
   \mathcal{W}_{ k + N  - \mathsf{rk}  \left(   v \right)  +1  }  \left(     F  \left[  X   \right]   \right)  
  \subseteq   \mathcal{W}_{ k + N  }  \left(     F  \left[  X   \right]   \right)  
  \      .
  \]
Since  $  F$ is an  $L_{  \mathbb{T} }^-  $-embedding, this means that we also have  
\begin{multline*}
 w_0    \in     \mathcal{W}_{ k + N  - \mathsf{rk}  \left(   v \right)  +1  }  \left(     F  \left[  X   \right]   \right)  
\subseteq  
   \mathcal{W}_{ k + N  - \mathsf{rk}  \left(   s \right)   }  \left(     F  \left[  X   \right]   \right)  
\     \mbox{  and   }      \   
\\
 w_1    \in     \mathcal{W}_{ k + N  - \mathsf{rk}  \left(   v \right)  +1  }  \left(     F  \left[  X   \right]   \right)  
\subseteq  
   \mathcal{W}_{ k + N  - \mathsf{rk}  \left(   t \right)   }  \left(     F  \left[  X   \right]   \right)  
 \mathcal{W}_{k + N }  \left( X  \right)    
   \end{multline*}
such that   
\[
   w  =  \langle  w_0, w_1  \rangle^{  \mathcal{B}  }  
\    \mbox{  and   }          \   
\langle  F(w_0) ,   F(w_1)  \rangle^{  \mathcal{B}  }  
= F(w)  =  L(v)  
= 
\langle  L(s), L(t)   \rangle^{  \mathcal{B}  } 
\     .
\]
Since  $  \langle  \cdot  ,  \cdot   \rangle^{  \mathcal{B}  }   $ is one-to-one,  from this  we get 
\[
 F(w_0)  =  L(s)      \    \mbox{  and  }      \         F(w_1)  = L(t)  
 \       .
 \]
By the induction hypothesis,   $  w =  s  $  and  $  w_1  =  t  $,  
and thus  
\[
  w  =  \langle  w_0,  w_1  \rangle^{  \mathcal{B} }   =     \langle  s,  t \rangle^{  \mathcal{B} }     =  v  
  \     .
  \]
 







We consider  cases    (2)  and  (3)  together since they are very similar. 
Let   
\[
\mathcal{G} :=   \begin{cases}
\mathcal{R}_m      &    \mbox{ in case of  }  (2) 
\\
\mathcal{A}_n      &    \mbox{ in case of  }  (3) 
\end{cases}
\]
Recall that    $  \mathcal{G}   $  is one-to-one. 
Since  $L$ is an $L_{  \mathbb{T} }^-  $-homomorphism  and  $F$  is  an   $L_{  \mathbb{T} }^-  $-embedding
\[
\mathcal{G}  \left(   L(s)  \right)     =   
 L(v)  = F(w)   \in  
 \mathcal{W}_{k + N- \mathsf{rk} \left(   v \right)  }  \left(   F  \left([ X  \right]   \right)   
 \     . 
\]
Since  $   k + N- \mathsf{rk}  \left(   v \right)  +1  \leq  k+N $  as  $   \mathsf{rk}  \left(   v \right)   >  0  $
\[
  L(s)  \in  
   \mathcal{W}_{k + N- \mathsf{rk} \left(   s \right)  }  \left(   F  \left([ X  \right]   \right)   
  \subseteq     \mathcal{W}_{ k + N  }  \left(     F  \left[  X   \right]   \right)  
  \      .
  \]
Since  $  F$ is an  $L_{  \mathbb{T} }^-  $-embedding, this means that we also have  
\[
u     \in   \mathcal{W}_{k + N -  \mathsf{rk} \left(   s \right)  }  \left( X  \right)    
 \   \mbox{ such that  }    \   
 w  =   \mathcal{G}   \left(  u  \right) 
   \      .
   \]
and   
\[
\mathcal{G}  \left(   F(u)  \right)     =   
 F(w)  =  L(v)    = 
\mathcal{G}  \left(   L(s)  \right)     
\         .
\]
Since  $   \mathcal{G} $ is one-to-one, $  F(u)  =  L(s)  $. 
By the induction hypothesis,   $u =  s $,     and thus  
\[
  w  =    \mathcal{G} \left(  u   \right)     =       \mathcal{G}   \left(  s   \right)     =   v  
  \     .
  \]




Thus,  by induction,   the claim holds. 



 
 
 
 We use the claim to prove the lemma. 
Pick    $  w,  v    \in   \mathsf{Cl}_{N}  \left(  \mathcal{W}_{k}  \left( X ,  \mathsf{d}^0_0  \right)   \right) $  such that  $  L(w)  =  L(v )  $. 
We prove by induction on  $  \mathsf{rk}  \left( w  \right)  +   \mathsf{rk}  \left( v  \right)   $  that  $  w  = v $. 
First,  we  consider the case  $  \min  \left(   \mathsf{rk}  \left( w  \right)     \,  ,  \,    \mathsf{rk}  \left( v  \right)     \right)  =  0   \,    $.
If    one of  $  w,  v  $  is an element of  $  \big\{  \mathsf{d}^r_0  :   \  r  \in  \mathbb{Z}  \,  \big\}  $, 
then the other must be an element of  $  \mathcal{W}_{k}  \left( X ,  \mathsf{d}^0_0  \right)   $,   
and we get  that     $  w  = v  $   since   $F$ is   \textsf{Good}. 
If  one of     $  w,  v  $   is   an element  of    $  \mathcal{W}_{k}  \left( X  \right)   $,  
then  $   w   =   v  $  by the claim  since  $F$ and $L$ agree on   $  \mathcal{W}_{k}  \left( X  \right)   $. 





We consider the case    $   \mathsf{rk}  \left( w  \right)    > 0  $  and     $    \mathsf{rk}  \left( v  \right)    >  0  $. 
We    have the following cases
\begin{enumerate}
\item    $  w  =  \langle  w_0 ,  w_1  \rangle^{  \mathcal{B}  }  $  and  $  v =  \langle  v_0, v_1  \rangle^{  \mathcal{B}  }  $
where  $ w_0, w_1,  v_0, v_1  \in     \mathsf{Cl}_{N}  \left(  \mathcal{W}_{k}  \left( X  \right)   \right)     $, 
 $  \mathsf{rk} \left(  w  \right)  =  1  +  \max\left(   \mathsf{rk} \left(  w_0 \right)  \,  ,  \,    \mathsf{rk} \left(  w_1  \right)   \right)     $ 
and  $  \mathsf{rk} \left(  v  \right)  =  1  +  \max\left(   \mathsf{rk} \left(  v_0 \right)  \,  ,  \,    \mathsf{rk} \left(  v_1  \right)   \right)    \,   $. 
  



\item     $  w  =    \mathcal{G}  \left(  g  \right)  $  and  $  v =  \langle v_0, v_1  \rangle^{  \mathcal{B}  }  $
where  $ g,  v_0, v_1  \in     \mathsf{Cl}_{N}  \left(  \mathcal{W}_{k}  \left( X  \right)   \right)     $
  $  \mathcal{G}     \in  \lbrace  \mathcal{R}_m   :  \   m  \in  \mathbb{Z} \,   \rbrace   \cup 
 \lbrace  \mathcal{A}_n :  \   n  \in    \omega  \setminus \lbrace 0, 1 ,  2   \rbrace  \,   \rbrace      $, 
  $  \mathsf{rk} \left(  w  \right)  =  1  +  \mathsf{rk} \left(  g \right)    $
and  $  \mathsf{rk} \left(  v  \right)  =  1  +  \max\left(   \mathsf{rk} \left(  v_0 \right)  \,  ,  \,    \mathsf{rk} \left(  v_1  \right)   \right)    \,   $. 



\item   $  w  =    \mathcal{G}  \left(  g   \right)  $  and  $  v =  \mathcal{H}  \left(  u \right)   $
where  $ u,  g  \in     \mathsf{Cl}_{N}  \left(  \mathcal{W}_{k}  \left( X  \right)   \right)     $
and  $  \mathcal{G} ,   \mathcal{H}      \in  \lbrace  \mathcal{R}_m   :  \   m  \in  \mathbb{Z} \,   \rbrace   \cup 
 \lbrace  \mathcal{A}_n :  \   n  \in    \omega  \setminus \lbrace 0, 1 ,  2    \rbrace  \,   \rbrace      $, 
   $  \mathsf{rk} \left(  w  \right)  =  1  +  \mathsf{rk} \left(  g \right)    $
and    $  \mathsf{rk} \left(  v  \right)  =  1  +  \mathsf{rk} \left(  u \right)       \,   $.
\end{enumerate} 
 We consider case (1). 
Since  $  L $  is an $  L_{  \mathbb{T} }^-  $  homomorphism,  we get  
\[
 \langle  L(w_0) ,  L( w_1)  \rangle^{  \mathcal{B}  }  
 =   L(w)   =  L(v)   =  
  \langle  L(v_0) ,  L( v_1)  \rangle^{  \mathcal{B}  }  
\     .
\]
Since  $  \langle  \cdot  ,  \cdot   \rangle^{  \mathcal{B}  }   $  is one-to-one,  
this implies  $  L(w_0) = L(v_0) $  and  $  L(w_1) = L(v_1) $. 
By the induction hypothesis,   we get $  w_0 = v_0 $  and  $  w_1= v_1 $, 
which implies  $  w = v $. 

 
 
 
 
 We consider case (2). 
 Let  
 \[
  \mathcal{G}^{  \prime  }    =   \begin{cases}
  \mathcal{R}_{ m-1}     &   \mbox{  if  }      \mathcal{G} =  \mathcal{R}_m    \mbox{  where  }   m  \in  \mathbb{Z}
  \\
    \mathcal{A}_{ n-1}     &   \mbox{  if  }      \mathcal{G} =  \mathcal{A}_n    \mbox{  where  }   n  \in  \omega  \setminus  \lbrace 0, 1,  2  \rbrace
\         .
  \end{cases}
 \]
 Since  $  L $  is an $  L_{  \mathbb{T} }^-  $  homomorphism,  we get  
\[
\langle    \mathcal{G}^{  \prime  }      \left(  L(g)  \right)    \,    ,    \,   L(g)   \rangle^{  \mathcal{B}  }  
= 
 \mathcal{G} \left(  L(g)  \right) 
 =   L(w)   =  L(v)   =  
  \langle  L(v_0) ,  L( v_1)  \rangle^{  \mathcal{B}  }  
\     .
\]
Since  $  \langle  \cdot  ,  \cdot   \rangle^{  \mathcal{B}  }   $  is one-to-one,  
this implies
\[
 \mathcal{G}^{  \prime  }      \left(  L(g)  \right)    =      L(v_0)   
 \     \mbox{  and   }       \    
    L(g)          =    L( v_1 )  
    \         .       \tag{*}
\]
Since    $  \mathcal{G}      \left(  g  \right)     \in     \mathsf{Cl}_{N}  \left(  \mathcal{W}_{k}  \left( X  ,  \mathsf{d}^0_0  \right)   \right)   $  
and $  \mathsf{rk}  \left(    \mathcal{G}      \left(  g   \right)   \right)   >  0  $, 
we also have  
 \[
   \mathcal{G}^{  \prime }      \left(   g \right)     \in     \mathsf{Cl}_{N}  \left(  \mathcal{W}_{k}  \left( X   ,  \mathsf{d}^0_0 \right)   \right)   
   \    \     \mbox{  and  }      \        \   
   \mathsf{rk}  \left(    \mathcal{G}      \left(  g \right)   \right)     \geq   
\mathsf{rk}  \left(    \mathcal{G}^{  \prime }       \left(  g  \right)   \right)   
\        .
\]
Hence,  from (*)   we get  
\[
  L   \left(    \mathcal{G}^{  \prime  }      \left(  g   \right)   \right)  
  =     \mathcal{G}^{  \prime  }      \left(  L(g)  \right)    =      L(v_0)   
 \     \mbox{  and   }       \    
    L(g)          =    L( v_1 )  
    \         .       
\]
By the induction hypothesis 
\[
  \mathcal{G}^{  \prime  }      \left(  g   \right)    =  v_0 
 \     \mbox{  and   }       \    
   g        =   v_1 
    \         .       
\]
Thus 
\[
w  =     \mathcal{G}     \left(  g   \right)  
= 
\langle    \mathcal{G}^{  \prime  }      \left( g   \right)    \,    ,    \,  g   \rangle^{  \mathcal{B}  }  
= 
\langle   v_0,  v_1   \rangle^{  \mathcal{B}  }  
=  v  
\          .
\]
 
 
 
 
 
 
 Finally,  we consider case (3). 
 Since  $ L$  is  an $L_{  \mathbb{T}  }^-$-homomorphism 
\[
 \mathcal{G}    \left(  L(g)  \right)      =  L(w)   =  L(v)   =     \mathcal{H}    \left(  L(u)  \right)  
\     .
\]
Since   any two distinct maps in   
$  \big\lbrace   \mathcal{R}_m   :   \   m  \in  \mathbb{Z}   \,   \rbrace    \cup  
 \big\lbrace   \mathcal{A}_n   :   \   n   \in     \omega  \setminus  \lbrace 0,  1  \rbrace   \,   \rbrace         $ 
 have disjoint images,   
 $   \mathcal{G}    =    \mathcal{H}  $. 
 This means  
 \[
  \mathcal{G}    \left(  L(g)  \right)   =   \mathcal{G}    \left(  L(u)  \right) 
  \      .
 \] 
 Since     $    \mathcal{G}    $  is one-to-one,   this implies  
 $  L(g)  =  L(u)  $. 
 By the induction hypothesis,  $  g  =  u  $, and thus 
 \[
 w  =    \mathcal{G}    \left(  g \right)   =   \mathcal{G}    \left(  u \right)     =   v 
 \       .
 \]
 
 
 
 
 
 
 
 
 
Thus,  by induction,  $L$ is one-to-one.
\end{proof}

\end{document}
